% Part: sets-functions-relations
% Chapter: functions
% Section: kinds
% Hindi translation (hi-Deva-IN) of Open Logic commit 9620cc73f9c8e0ad003c514a5d3748f29611c4c0.

\documentclass[../../../include/open-logic-section]{subfiles}

\begin{document}

\olfileid[hi]{sfr}{fun}{kin}
\olsection{फलनों के प्रकार}

\begin{explain}
जिन फलनों के प्रकारों से हमारा सबसे अधिक सामना होता है, उनमें से कुछ का
वर्गीकरण प्रस्तुत करना उपयोगी होगा।

आरंभ में हम ऐसे फलनों पर विचार कर सकते हैं जिनमें सहप्रांत का प्रत्येक अवयव
फलन का कोई मान होता है। ऐसे फलन !!{surjective} कहलाते हैं और उन्हें
\olref{fig:surjective} की तरह चित्रित किया जा सकता है।

\begin{figure}[H]
  \olasset{assets/diagrams/surjective.tikz}
  \caption{!!^a{surjective} फलन में सहप्रांत का प्रत्येक !!{element} कोई
    मान होता है।}
  \ollabel{fig:surjective}
\end{figure}
\end{explain}

\begin{defn}[!!^{surjective} फलन]
फलन $f \colon A \rightarrow B$ \emph{!!{surjective}} है यदि और केवल यदि $B$
भी~$f$ का परिसर है, अर्थात् प्रत्येक $y \in B$ के लिए कम-से-कम एक
$x \in A$ ऐसा है कि~$f(x) = y$; प्रतीकों में:
\[
  (\forall y \in B)(\exists x \in A)f(x) = y.
\]
ऐसे फलन को $A$ से $B$ तक का !!a{surjection} कहते हैं।
\end{defn}

\begin{explain}
यदि आप यह दिखाना चाहते हैं कि $f$ !!a{surjection} है, तो आपको दिखाना होगा
कि~$f$ के सहप्रांत की प्रत्येक वस्तु $f(x)$ का मान है, किसी निवेश $x$ के लिए।

ध्यान दें कि प्रत्येक फलन एक !!a{surjection} \emph{प्रेरित करता है}। वास्तव
में, फलन $f \colon A \to B$ दिए जाने पर $f' \colon A \to \ran{f}$ को
$f'(x) = f(x)$ द्वारा परिभाषित करें। चूँकि $\ran{f}$ को
$\Setabs{f(x) \in B}{x \in A}$ के रूप में \emph{परिभाषित} किया गया है,
इसलिए यह फलन $f'$ निश्चित रूप से !!a{surjection} है।
\end{explain}

\begin{explain}
हर फलन प्रत्येक संभावित निवेश को एक अद्वितीय निर्गम पर प्रतिचित्रित करता है।
पर ऐसे फलन भी होते हैं जो अलग-अलग निवेशों को कभी एक ही निर्गम पर प्रतिचित्रित
नहीं करते। ऐसे फलन !!{injective} कहलाते हैं और उन्हें
\olref{fig:injective} की तरह चित्रित किया जा सकता है।
\begin{figure}[H]
  \olasset{assets/diagrams/injective.tikz}
  \caption{!!^a{injective} फलन दो अलग-अलग फलन-कोणांकों को कभी एक ही मान पर
    प्रतिचित्रित नहीं करता।}
  \ollabel{fig:injective}
\end{figure}
\end{explain}

\begin{defn}[!!^{injective} फलन]
फलन $f \colon A \rightarrow B$ \emph{!!{injective}} है यदि और केवल यदि
प्रत्येक $y \in B$ के लिए अधिक-से-अधिक एक $x \in A$ ऐसा है कि~$f(x) = y$।
ऐसे फलन को $A$ से~$B$ तक का !!a{injection} कहते हैं।
\end{defn}

\begin{explain}
यदि आप यह दिखाना चाहते हैं कि $f$ !!a{injection} है, तो आपको दिखाना होगा कि
किन्हीं !!{element}ों $x$ और $y$ के लिए, जो $f$ के प्रांत में हों, यदि
$f(x)=f(y)$, तो $x=y$।
\end{explain}

\begin{ex}
$f\colon \Nat \to \Nat$ अचर फलन, जो $f(x) = 1$ द्वारा दिया गया है, न तो
!!{injective} है, न !!{surjective}।

$f\colon \Nat \to \Nat$ तत्समक फलन, जो $f(x) = x$ द्वारा दिया गया है,
!!{injective} और !!{surjective} दोनों है।

$f \colon \Nat \to \Nat$ परवर्ती फलन, जो $f(x) = x+1$ द्वारा दिया गया है,
!!{injective} है, पर !!{surjective} नहीं।

निम्न प्रकार से परिभाषित फलन $f \colon \Nat \to \Nat$:
\[
  f(x) =
  \begin{cases}
    \frac{x}{2} & \text{यदि $x$ सम है} \\
    \frac{x+1}{2} & \text{यदि $x$ विषम है।}
  \end{cases}
\]
!!{surjective} है, पर !!{injective} नहीं।
\end{ex}

\begin{explain}
प्रायः हम ऐसे फलनों पर विचार करना चाहते हैं जो !!{injective} और
!!{surjective} दोनों हों। ऐसे फलनों को हम !!{bijective} कहते हैं। वे
\olref{fig:bijective} में चित्रित फलन जैसे दिखते हैं। !!^{bijection}s ऐसे
प्रतिचित्रण हैं जिन्हें कभी-कभी \emph{एकैकी संगतियाँ} भी कहते हैं, क्योंकि वे सहप्रांत के अवयवों को
प्रांत के अवयवों के साथ अद्वितीय रूप से युग्मित करते हैं।
\begin{figure}[H]
  \olasset{assets/diagrams/bijective.tikz}
  \caption{!!^a{bijective} फलन सहप्रांत के अवयवों को प्रांत के अवयवों के साथ
    अद्वितीय रूप से युग्मित करता है।}
  \ollabel{fig:bijective}
\end{figure}
\end{explain}

\begin{defn}[!!^{bijection}]
फलन $f \colon A \to B$ \emph{!!{bijective}} है यदि और केवल यदि वह
!!{surjective} और !!{injective} दोनों है। ऐसे फलन को $A$ से~$B$ तक (या
$A$ और~$B$ के बीच) !!a{bijection} कहते हैं।
\end{defn}

\end{document}
